% ============================================================
%  Section 2, Segment 2B — Chords \& Chord Progressions / 和弦与进行
%  原书页码 27-53
% ============================================================
\chapter{Chords and Chord Progressions / 和弦与和弦进行}
\label{ch:chords}

\eng{By now, you should have an understanding of how basic chords are built and how they function in major and minor keys. The next step is to take these basic chords and make them into extended chords. Extended chords are used a ton in mathy and emo styles of music. They add a bit of extra flair to basic chords, so they have the potential to play with emotions and feelings in different ways.}

\chn{到现在，你应该已经理解了基本和弦是如何构建的，以及它们在大调和小调中的功能。下一步是把这些基本和弦扩展成延伸和弦。延伸和弦在数学摇滚和情绪摇滚风格中大量使用。它们为基本和弦增加了一点额外的色彩，因此有能力以不同方式影响情感与氛围。}

\paragraph{Extended Chords}
\eng{Extended chords are made by adding extra notes to basic chord spellings. Common extensions begin with the seventh chords. Again, like basic major, minor and diminished chords, extended chords have their own spellings. Below are the chord spellings for extended seven chords, the first possible extension:}

\chn{延伸和弦是通过在基本和弦构成上添加额外音符来形成的。常见的扩展从七和弦开始。同样，像基本的大三和弦、小三和弦和减和弦一样，延伸和弦也有自己的构成方式。下面是延伸七和弦的和弦记号，这是第一个可能的扩展：}

\begin{itemize}[leftmargin=*]
  \item \eng{Major 7th: 1–3–5–7}\chn{大七和弦：1–3–5–7}
  \item \eng{Dominant 7th: 1–3–5–b7}\chn{属七和弦：1–3–5–b7}
  \item \eng{Minor 7th: 1–b3–5–b7}\chn{小七和弦：1–b3–5–b7}
  \item \eng{Minor 7 flat 5: 1–b3–b5–b7}\chn{半减七和弦：1–b3–b5–b7}
\end{itemize}

\eng{You may have noticed that we have added two new chord types: dominant seven, and a minor seven flat five chord. Dominant chords have a strong sound to them and feel they need to resolve to another chord. M7b5 chords add more suspense to the basic diminished chord. They are useful as tension builders, and for slotting in-between chords in chord progressions for some suspenseful colouring.}

\chn{你可能已经注意到，我们增加了两种新的和弦类型：属七和弦，以及半减七和弦。属七和弦有一种强烈的声音，给人一种需要解决到另一个和弦的感觉。m7b5 和弦为基本的减和弦添加了更多悬念。它们可作为张力构建器，并用于和弦进行中插入在两个和弦之间，以获得悬念感。}

\eng{Beyond extended sevens, there are extended nines, elevens, and even thirteenth chords! Realistically it is up to you to choose which extended chords you want to learn and use. I recommend learning all the shapes provided in this guide and over time find which work for you. For example, I found I largely prefer to use extended major and minor nines, and minor seven and minor eleven chords.}

\chn{在延伸七和弦之外，还有延伸九和弦、十一和弦，甚至十三和弦！实际上由你决定想学习和使用哪些延伸和弦。我建议学习本指南中提供的所有形状，并随着时间推移发现哪些最适合你。例如，我发现我主要喜欢使用延伸的大九、小九，以及小七和小十一和弦。}

\paragraph{Why Are the 9ths, 11ths, and 13ths Called That?}
\eng{Before we continue, you may be saying to yourself “hang on, there are only seven notes in a scale, so where are the 9ths, 11ths, 13ths coming from?” Let me elaborate. Counting through the scale of C major we eventually arrive at the octave (the same note but higher in pitch). If we continue counting past the octave, notes 2, 4, or 6 are referred to as 9th, 11th, and 13th degrees when used in chords. This is because they are extended above the octave.}

\chn{在继续之前，你可能会问“等一下，音阶里只有七个音，九度、十一度、十三度是从哪里来的？”让我来解释。沿着 C 大调音阶数下去，我们最终会到达八度（同一个音，但更高的音高）。如果继续数过八度，第二个、第四个或第六个音在和弦中使用时会被称为九度、十一度和十三度。这是因为它们是高于八度的延伸音。}

\eng{It is important to note that extended chords like 9, 11, and 13 are named that way only if they include the seventh degree. For instance, if you have a chord with the notes 1, 3, 5, and 9, it is simply called an Add9 chord, like Cmajadd9 for C major. The seventh degree is crucial for the sound of these extended chords.}

\chn{重要的是要注意，只有当和弦包含七度时，九度、十一度、十三度这样的延伸和弦才会这么命名。例如，如果你有一个包含 1、3、5 和 9 的和弦，它只叫做 Add9 和弦，就像 Cmajadd9。七度对于这些延伸和弦的音色至关重要。}

\eng{The reason I mention this is that with larger extended chords, you have more notes, but limited fingers and fretboard space to play them all. So, some “less important” notes are often left out to make it playable. However, it is noteworthy that the seventh degree is usually retained in these chords.}

\chn{我之所以提到这个，是因为在更大的延伸和弦中，你有更多的音，但手指和指板空间有限，无法全部弹出。因此，为了可演奏性，有些“次要”的音常常会被省略。然而值得注意的是，七度通常会保留在这些和弦中。}

\paragraph{Common Extended Chord Notes}
\eng{The following table shows how extended chords are often simplified for playability, with certain notes omitted while retaining the seventh degree.}

\chn{下面的表格显示了延伸和弦如何为了便于演奏而简化，省略某些音但保留七度。}

\begin{itemize}[leftmargin=*]
  \item \eng{Major 7: 1–3–5–7}\chn{大七：1–3–5–7}
  \item \eng{Major 9: 1–3–5–7–9, often dropping the 5th to become 1–3–7–9}\chn{大九：1–3–5–7–9，通常省略五度变成 1–3–7–9}
  \item \eng{Minor 11: 1–b3–5–b7–9–11, often dropping the 5th and 9th to become 1–b3–b7–11}\chn{小十一：1–b3–5–b7–9–11，通常省略五度和九度变成 1–b3–b7–11}
  \item \eng{Dom 13: 1–3–5–b7–9–11–13, often dropping the 5th and 11th (sometimes the 9th) resulting in 1–3–b7–(9)–13}\chn{属十三：1–3–5–b7–9–11–13，通常省略五度和十一度（有时也省略九度），结果为 1–3–b7–(9)–13}
\end{itemize}

\paragraph{Seventh Chord Voicings \& Practice Progressions}
\eng{Below are extended seven chord voicings starting from the sixth, fifth, and fourth strings. Learning in all three positions will help with changing chords in progressions more smoothly and provide you with more options. So take your time to learn them, and keep in mind there are other ways of voicing these chords. These are shapes I like, and I think you will too.}

\chn{下面是从第六、第五月和第四弦开始的延伸七和弦指型。学习这三种位置会帮助你更顺畅地在和弦进行中转换，并为你提供更多选项。所以慢慢学习它们，并记住还有其他方式来配位这些和弦。这些是我喜欢的形状，我想你也会喜欢。}

\eng{In common music practice, changes between chords in a progression are usually accomplished by moving to the closest available voicing of a particular chord. There are a couple of reasons for this: smoothness of overall feel and less effort for the guitar player. However, interesting things can be achieved with bigger chord jumps, and this happens quite often in math rock. So do not take this as a rule that must always be followed.}

\chn{在常见的音乐实践中，和弦进行中的变换通常通过移动到某个和弦的最近配位来完成。这有两个原因：整体感觉更平滑，以及吉他手的努力更少。然而，通过更大的和弦跳跃也可以实现有趣的效果，这在数学摇滚中经常发生。所以不要把这当作必须始终遵循的规则。}

\eng{Here is a chord progression in the key of C major (CM7, Am7, G7, FM7) represented as I, vi, V, IV. I have shown three ways to play this progression on the guitar neck, starting from different initial chords. Try playing them using your preferred strumming patterns. Additionally, experiment with other ways to play the progression, possibly with larger chord jumps or spacing between the chords.}

\chn{这是一个 C 大调中的和弦进行（CM7、Am7、G7、FM7），表示为 I、vi、V、IV。我展示了三种不同的方式在吉他指板上演奏这一进行，起始和弦不同。尝试用你喜欢的扫弦模式演奏它们。还可以尝试以其他方式演奏这个进行，比如更大的和弦跳跃或在和弦之间留出更多间距。}

\paragraph{Ninth Chord Voicings \& Practice Progressions}
\eng{Extended nine chords are hands down my favourite chord type to use. They have something about them that makes them just a bit more special than extended seventh chords. That being said, I almost always use a mixture of chord types in my chord progressions, because some chord shapes or chord types work better in a particular place in a progression than others.}

\chn{延伸九和弦无疑是我最喜欢使用的和弦类型。它们有一种使它们比延伸七和弦更特别的感觉。话虽如此，我几乎总是在和弦进行中混合使用多种和弦类型，因为某些和弦形状或类型在进行的特定位置上比其他更合适。}

\eng{Ninth chords have the spelling 1, 3, 5, 7, 9 and like extended seventh chords, I have included the spellings for each variation for clarity:}

\chn{九和弦的构成是 1、3、5、7、9。和延伸七和弦一样，为了清楚起见，我也列出了每种变化的构成：}

\begin{itemize}[leftmargin=*]
  \item \eng{Major 9th: 1–3–5–7–9}\chn{大九：1–3–5–7–9}
  \item \eng{Dominant 9th: 1–3–5–b7–9}\chn{属九：1–3–5–b7–9}
  \item \eng{Minor 9th: 1–b3–5–b7–9}\chn{小九：1–b3–5–b7–9}
  \item \eng{Minor 9 flat 5: 1–b3–b5–b7–9}\chn{小九降五：1–b3–b5–b7–9}
\end{itemize}

\eng{As you can see, the ninth degree of the chord is not altered (sharp or flat) in any of the extended ninth chords. The same is true for extended elevens and thirteenth chords. This is because the guide tones of a chord are its third and seventh. Even without a root you can almost completely define a chord by just knowing its third and seventh.}

\chn{如你所见，在所有延伸九和弦中，第九度并没有被改变（升号或降号）。对于延伸十一和弦和十三和弦也是如此。这是因为和弦的导音是三度和七度。即使没有根音，仅凭三度和七度你也几乎可以完全定义一个和弦。}

\eng{Note: The minor nine flat five is a little different because of its flattened fifth. It still contains the same guide notes as a minor nine chord: a flattened third and seventh, but until we add the fifth (either a natural fifth or a flattened fifth) into the chord we will not know for sure if the chord is a minor or a half-diminished chord (m7b5).}

\chn{注意：小九降五与众不同，因为它有降五度。它仍然包含与小九和弦相同的导音：降三度和七度，但是在加入五度（无论是纯五度还是降五度）之前，我们无法确定该和弦究竟是小和弦还是半减和弦（m7b5）。}

\eng{Another useful application of these chords is to substitute basic major, minor, or diminished chords or to use in conjunction with them for a more colourful sound.}

\chn{这些和弦的另一个有用应用是替换基本的大三和弦、小三和弦或减和弦，或者与它们结合使用，以获得更丰富的音色。}

\paragraph{Other Chord Variations Useful for Math Rock}
\eng{There are many more chords and chord types to explore beyond what we have discussed. Even without delving into altered chords and chord inversions, there is still a vast world of chords to discover. In this section, I will highlight some chords that I find particularly useful in mathy styles of music, and I will provide guidance on when and where to incorporate them.}

\chn{除了我们已经讨论的内容之外，还有许多更多的和弦和和弦类型可以探索。即使不深入改变和弦和和弦转位，仍然有一个广阔的和弦世界可以发现。在本节中，我将重点介绍一些我认为在数学摇滚风格中尤其有用的和弦，并提供何时何地使用它们的建议。}

\paragraph{Suspended Chords}
\eng{A suspended chord is a chord where the third degree has been omitted and usually replaced with either a major second or a perfect fourth. They have neither a major nor minor sound to them; rather, more of an “open sound” that often feels like it needs to be resolved.}

\chn{挂留和弦是省略三度并通常用大二度或纯四度代替的和弦。它们既没有大调也没有小调的音色，而是一种更“开放”的声音，通常给人一种需要解决的感觉。}

\eng{The most popular usage of suspended chords in math rock and Midwest emo is to embellish an existing chord in a progression, creating movement within the chord itself, or use them to situate the listener in a feeling of uncertainty. For example, if you were to play an A major chord followed by an Asus4 you will hear how this embellishes the original A major chord, giving it a different feel and flavour. In addition, you can play this backwards (Asus4 to A) for a more resolved sound.}

\chn{挂留和弦在数学摇滚和中西部情绪摇滚中最受欢迎的用法是装饰和弦进行中的现有和弦，创造和弦内部的运动，或者将听众置于一种不确定的感觉中。例如，如果你先弹 A 大调和弦再弹 Asus4，你会听到它如何装饰原来的 A 大调，赋予它不同的感觉与色彩。此外，你也可以把它反过来弹（Asus4 到 A）以获得更解决感的声音。}

\eng{You could even add an Asus2 to the progression. Listen how the slight change of the ascending higher note alters the overall feel of the chord.}

\chn{你甚至可以在进行中加入 Asus2。听一听升高音的微小变化如何改变和弦的整体感觉。}

\eng{Another useful application of these chords is to use them to replace dominant chords that sound out of place in a chord progression or to replace an unwanted dominant sound in a progression. Remember, dominant seven chords are built from the fifth degree (V chord) of a key. Try the progressions below to hear this in action. Personally, I find the dominant chord sounds a little out of place for mathy sounds, but the suspended four chord replacement makes the progressions much more suitable for the desired sound I was looking for.}

\chn{这些和弦的另一个有用应用是用它们替换在和弦进行中听起来不合适的属和弦，或者替换进行中不想要的属和弦声音。记住，属七和弦是从一个调的五度（V 和弦）构建的。尝试听下面的进行来感受这一点。就个人而言，我发现属和弦对数学摇滚的声音来说有点不合适，但用挂留四和弦替换后，使这些进行更适合我想要的音色。}

\eng{One last point on suspended chords. It is worth pointing out that not every triad chord in a key can be replaced with either a sus2 or sus4 chord. I have made a handy chart below so you can see which chords can be substituted with a sus chord.}

\chn{关于挂留和弦的最后一点值得指出：并不是调中的每个三和弦都可以被 sus2 或 sus4 替代。我做了一个实用的图表，让你可以看到哪些和弦可以用 sus 和弦替代。}

\paragraph{Major and Minor 6 and 6add9 Chords}
\eng{Major and minor 6, add9, or 6add9 chords offer creative alternatives to extended chords while avoiding the seventh degree. They add harmonic interest without the complexity of including the seventh degree. Major 6 and 6add9 chords are often used because they are a great substitute for major seventh chords when it comes to providing a more resolved stable sound.} 

\chn{大/小六和弦、add9 和 6add9 和弦提供了比延伸和弦更有创意的替代方案，同时避免了七度。它们在不引入七度复杂性的情况下增加了和声趣味。大六和 6add9 和弦经常使用，因为它们在提供更解决、稳定的声音方面，是大七和弦的极佳替代。}

\eng{The reason is that a major seventh chord contains a minor second interval between the root and the major seventh degree. Minor second intervals do not sound pleasant. If you play a major seventh chord, the interval is still there, but it is more spaced out and less obvious when the other notes are added. In contrast, major 6 chords and major 6add9 chords avoid this close interval and still produce a colourfully rich sound.}

\chn{原因在于大七和弦包含根音与大七度之间的一个小二度音程。小二度听起来并不悦耳。如果你弹一个大七和弦，这个音程依然存在，但它更分散，并且在加入其他音符时不那么明显。相比之下，大六和弦和大六加九和弦避免了这个紧密音程，同时仍然产生丰富的色彩声音。}

\eng{It is worth keeping in mind that minor 6 and minor 6add9 chords can have a more dissonant sound compared to their minor 7 counterparts. Therefore, it is a good idea to experiment with them to achieve the specific musical results you desire. Another useful application of these chords is to substitute basic major, minor, or diminished chords or to use them in conjunction with other chord types for a more colourful sound.}

\chn{值得记住的是，与小七和弦相比，小六和小六加九和弦可能具有更加不协和的声音。因此，尝试它们是实现你想要的特定音乐效果的好办法。这些和弦的另一个有用应用是替换基本的大三、小三或减和弦，或者与其他和弦类型结合使用，以获得更丰富的色彩。}

\paragraph{Extended Shell Chord Voicings}
\eng{Extended shell chord voicings simplify extended chords by including only the essential notes: the root, third, and seventh. These voicings capture the core tonal quality of the chord while removing some of the extra notes. They are particularly useful in ensemble settings with complex harmony and also work well with distortion, making them great for chord-based riffs.}

\chn{延伸壳和弦配位通过仅包含必要的音：根、三度和七度，来简化延伸和弦。这些配位保留了和弦的核心音色，同时省略了一些额外音符。它们在复杂和声的合奏环境中尤其有用，也适合失真声音，是和弦型 riff 的理想选择。}

\eng{Major seven shell voicings, minor seven shell voicings, and dominant seven shell voicings are all useful shapes to learn. Minor seven flat five shell voicings are the same as minor seven shell voicings, since the fifth is omitted in the shell voicing construction. This is not a problem; the shell voicing still functions as a seventh chord in the harmonic context.}

\chn{大七壳配位、小七壳配位和属七壳配位都是值得学习的形状。半减七壳配位与小七壳配位相同，因为在壳配位构成中省略了五度。这不是问题；在和声语境中，壳配位仍然可以作为七和弦发挥作用。}

\eng{Shell voicings are worth practicing alongside regular full-sized chords. Use them as a basis for crafting riffs or as a source of chord-writing ideas. They have a strong and emotive quality, especially with distortion or gain.}

\chn{壳配位值得与常规全尺寸和弦一起练习。把它们用作创作 riff 的基础，或作为和弦写作的灵感来源。它们具有强烈而富有表现力的特质，尤其是在失真或增益下。}

\paragraph{Chord Progressions \& Mixing Chord Types}
\eng{A chord progression is a sequence of chords that forms the backbone of a song section. Some progressions sound more pleasing than others, and they can range from simple to highly intricate. Knowing the key helps you choose chords that will fit together well.}

\chn{和弦进行是一系列和弦，构成歌曲片段的骨架。有些进行听起来比其他更悦耳，它们可以从简单到复杂不等。了解调有助于你选择能够和谐搭配的和弦。}

\eng{For example, in the key of C major, the diatonic chords are: C, Dm, Em, F, G, Am, Bdim. These correspond to I, ii, iii, IV, V, vi, vii°. In a major key, the I chord has the strongest pull because it resolves the key center, while the IV chord is stable but tends to move to another chord. The vi chord is the relative minor and often has a strong emotional pull.} 

\chn{例如，在 C 大调中，调内和弦是：C、Dm、Em、F、G、Am、Bdim。它们对应 I、ii、iii、IV、V、vi、vii°。在大调中，I 和弦具有最强的解决力，因为它会回到调中心，而 IV 和弦则稳定但通常倾向于移动到另一个和弦。vi 和弦是关系小调，通常具有强烈的情感吸引力。}

\eng{Below are several useful chord progressions based on common popular mathy and emo songs. These progressions are diatonic, meaning they only use chords from one key. They are not exact transcriptions of any particular song, but they capture useful harmonic functions.}

\chn{下面是一些基于常见数学摇滚和情绪摇滚歌曲的实用和弦进行。这些进行都是调内的，也就是说仅使用一个调内的和弦。它们并非任何特定歌曲的精确转录，但它们抓住了有用的和声功能。}

\paragraph{Functional Popular Diatonic Chord Progressions}
\eng{1. IV – V – vi – I. This progression has a satisfying ascending feeling and resolves nicely to the tonic. It works well for an ending, a middle section, or a chorus.}

\chn{1. IV – V – vi – I。这个进行有令人满意的上行感觉，最后很好地解决到主音。它适合作为结尾、中段或副歌。}

\eng{2. IV – ii – I – ii. This progression feels poppy with a touch of melancholy. It works well for a chorus-like section.}

\chn{2. IV – ii – I – ii。这个进行听起来流行但带有一丝忧郁。它适合作为副歌式的段落。}

\eng{3. IV – V – vi – V. This progression is useful for building tension and feels like it needs to move somewhere. The vi chord often acts as an emotional anchor.}

\chn{3. IV – V – vi – V。这个进行适合用来构建张力，听起来像是需要去往某处。vi 和弦通常起到情感锚点的作用。}

\eng{4. IV – ii – I. This progression works nicely for a verse or bridge, with a pleasant and stable feeling.}

\chn{4. IV – ii – I。这个进行很适合作为主歌或桥段，带来愉悦且稳定的感觉。}

\eng{5. I – IV – vi. This progression has an uplifting ascending feel.}

\chn{5. I – IV – vi。这个进行具有令人振奋的上行感觉。}

\eng{6. iii – IV – vi – V and 7. IV – V – iii – vi. These progressions have a math rock meets neo-soul vibe and work well with minor pentatonic or additional melodic colouring.}

\chn{6. iii – IV – vi – V 和 7. IV – V – iii – vi。这些进行具有数学摇滚与新灵魂的混合风格，搭配小调五声音阶或额外的旋律色彩效果很好。}

\eng{8. ii – IV – vi – I. This progression has a suspenseful math-pop feel.}

\chn{8. ii – IV – vi – I。这个进行具有悬念感的数学流行风格。}

\eng{9. I – iii – vi – ii. This is a nice progression for an epic outro.}

\chn{9. I – iii – vi – ii。这个进行适合作为史诗般的结尾。}

\eng{As a bonus, take a progression from Leon Russell's “Hummingbird” and adapt it for a minor key. The chord sequence i – VI – i – iv and i – V – VI – V has a bluesy original feel but can work well in a math rock context.}

\chn{作为奖励，可以参考 Leon Russell 的《Hummingbird》并把它改编为小调。和弦序列 i – VI – i – iv 和 i – V – VI – V 原本有布鲁斯风格，但在数学摇滚语境中也非常适用。}

\paragraph{Chord Progressions in Context (Standard Tuning)}
\eng{The following examples are derived from math rock and Midwest emo style songs. They are not exact reproduction, but they illustrate how many of the chords and progressions we discussed are used in real music.}

\chn{以下示例来自数学摇滚和中西部情绪摇滚风格的歌曲。它们不是精确再现，但展示了我们讨论的许多和弦与进行如何在实际音乐中使用。}

\eng{Midwest Emo, Progressive Math Rock, Japanese Math Rock, Emo Math Rock, Post-Hardcore/Math Rock all have their own characteristic progressions. Pay attention to the way the chords flow and the use of extended harmony.}

\chn{中西部情绪摇滚、渐进数学摇滚、日本数学摇滚、情绪数学摇滚、后硬核/数学摇滚都有各自的特点。注意和弦的流动方式以及延伸和声的使用。}

\paragraph{Chord Progressions Containing Borrowed Chords}
\eng{Borrowed chords are chords taken from another key than the parent key. They are often borrowed from the parallel minor or major. For example, if we are in C major, we could borrow chords from C minor. Borrowed chords add colour and unexpected movement to progressions.}

\chn{借用和弦是从主调之外的另一个调借用的和弦。它们通常来源于平行小调或大调。例如，如果我们在 C 大调中，可以借用 C 小调的和弦。借用和弦为进行增添了色彩和意料之外的变化。}

\eng{In the key of C major, the diatonic chords are I, ii, iii, IV, V, vi, vii°. In C minor, the chords are i, ii°, III, iv, v, VI, VII. Borrowing a chord from C minor when in C major introduces bIII, bVI, or bVII, which bring a darker or more dramatic character to the progression.}

\chn{在 C 大调中，调内和弦是 I、ii、iii、IV、V、vi、vii°。在 C 小调中，和弦是 i、ii°、III、iv、v、VI、VII。当在 C 大调中借用 C 小调的和弦时，会引入 bIII、bVI 或 bVII，它们会为进行带来更暗沉或更戏剧化的特质。}

\eng{Example borrowed progressions include:}

\begin{itemize}[leftmargin=*]
  \item \eng{IV – I – ii – bIII: the borrowed major III adds an unexpected colour.}\chn{IV – I – ii – bIII：借用的大三和弦增加了意外的色彩。}
  \item \eng{I – ii – IV – bVII: the borrowed bVII replaces the diminished feel with a more stable major chord.}\chn{I – ii – IV – bVII：借用的 bVII 用更稳定的大三和弦替代了减和弦的感觉。}
  \item \eng{IV – I – iii – v: borrowing the minor v gives the progression a more characterful sound.}\chn{IV – I – iii – v：借用的小五和弦为进行增添了更有特色的声音。}
  \item \eng{vi – I – ii – v: another variation with borrowed v.}\chn{vi – I – ii – v：另一种带借用 v 的变体。}
  \item \eng{I – vi – bVI – ii: the borrowed bVI leads smoothly back into the progression.}\chn{I – vi – bVI – ii：借用的 bVI 让进行顺畅地回到原路。}
\end{itemize}

\paragraph{Tasks Answers}
\eng{Here are the answers for the chord key tasks:}

\chn{以下是和弦调任务的答案：}

\begin{enumerate}[leftmargin=*]
  \item \eng{E major – A major – G\# minor – F\# minor – E major} \chn{E 大调 – A 大调 – G\# 小调 – F\# 小调 – E 大调} = \eng{E major}\chn{E 大调}
  \item \eng{D minor – G7 – C major} \chn{D 小调 – G7 – C 大调} = \eng{C major}\chn{C 大调}
  \item \eng{D major – E minor – F\# diminished – G major} \chn{D 大调 – E 小调 – F\# 减和弦 – G 大调} = \eng{G major}\chn{G 大调}
  \item \eng{C major – G major – A minor – E minor} \chn{C 大调 – G 大调 – A 小调 – E 小调} = \eng{Either C major or G major (I–V–vi–iii in C or IV–V–ii–vi in G)}\chn{可以是 C 大调，也可以是 G 大调（在 C 中为 I–V–vi–iii，在 G 中为 IV–V–ii–vi）}
  \item \eng{G major – B minor – C major – C minor} \chn{G 大调 – B 小调 – C 大调 – C 小调} = \eng{G major with a borrowed iv chord (I–iii–IV–iv)}\chn{带借用 iv 的 G 大调（I–iii–IV–iv）}
\end{enumerate}
